% ===== 原文 p.31 =====
% 译者 tr-04：表 4.7、表 4.8、例 4.9、§4.2.1 中 (B,A,A)-膜与正则纤维相交的概述
% 边界：本片末句补完 p31→p32 跨页句 "…Prym varieties of quotient / spectral curves, as shown in [89,91]."；
%       tr-05 从 p32 第 2 句 "Having mentioned the spectral data …" 起，不得重复。

% [原文 p.31]
\begin{table}[htbp]
\centering
\small
\begin{tabular}{llllll}
\toprule
$\mathfrak{g}_\C$ & $G_\C$ & 分裂形式 & 紧形式 $\mathfrak{u}$ & 固定 $\mathfrak{u}$ 的反对合 $\rho$ & $\dim\mathfrak{u}$ \\
\midrule
$\mathfrak{a}_n$ & & $\mathfrak{sl}(n,\R)$ & $\mathfrak{su}(n)$ & $\rho(X)=-\bar{X}^t$ & $n(n+1)$ \\
$\mathfrak{b}_n$ & & $\mathfrak{so}(n,n+1)$ & $\mathfrak{so}(2n+1)$ & $\rho(X)=\bar{X}$ & $n(2n+1)$ \\
$\mathfrak{c}_n$ & & $\mathfrak{sp}(2n,\R)$ & $\mathfrak{sp}(n)$ & $\rho(X)=\mathrm{J}_n\bar{X}\mathrm{J}_n^{-1}$ & $n(2n+1)$ \\
$\mathfrak{d}_n$ & & $\mathfrak{so}(n,n)$ & $\mathfrak{so}(2n)$ & $\rho(X)=\bar{X}$ & $n(2n-1)$ \\
\bottomrule
\end{tabular}
\tabtitle{4.7}{经典李代数的紧形式 $\mathfrak{u}$}
\end{table}

% [Tutor Note: p.31 表 4.7 逐格核对（列序：$\mathfrak{g}_\C$ | $G_\C$ | 分裂形式 | 紧形式 $\mathfrak{u}$ | 固定 $\mathfrak{u}$ 的反对合 $\rho$ | $\dim\mathfrak{u}$）：$\mathfrak{b}_n$ 行 $\dim\mathfrak{so}(2n+1)=n(2n+1)$、$\mathfrak{c}_n$ 行 $\dim\mathfrak{sp}(n)=n(2n+1)$、$\mathfrak{d}_n$ 行 $\dim\mathfrak{so}(2n)=n(2n-1)$ 均核对相符；但 $\mathfrak{a}_n$ 行的 $n(n+1)$ 与 $\dim\mathfrak{su}(n)=n^2-1$ 不符（疑应为 $n^2-1$），照原样排。]

于是，下表给出 $\theta$ 用复李代数上的反全纯对合表示时一切可能的取值。

\begin{table}[htbp]
\centering
\small
\begin{tabular}{llll}
\toprule
$\mathfrak{g}_\C$ & 实形式 $G$ & 固定 $G$ 的 $\sigma$ & $\theta=\rho\sigma$ \\
\midrule
\multirow{3}{*}{$\mathfrak{a}_n$} & $\Sl(n,\R)$ & $\sigma(X)=\bar{X}$ & $\theta(X)=-X^t$ \\
 & $\SU^*(2m)$ & $\sigma(X)=\mathrm{J}_m\bar{X}\mathrm{J}_m^{-1}$ & $\theta(X)=-\mathrm{J}_m X^t\mathrm{J}_m^{-1}$ \\
 & $\SU(p,q)$ & $\sigma(X)=-\mathrm{I}_{p,q}\bar{X}^t\mathrm{I}_{p,q}$ & $\theta(X)=\mathrm{I}_{p,q}X\mathrm{I}_{p,q}$ \\
\midrule
$\mathfrak{b}_n$ & $\SO(p,q)$ & $\sigma(X)=\mathrm{I}_{p,q}\bar{X}\mathrm{I}_{p,q}.$ & $\theta(X)=\mathrm{I}_{p,q}X\mathrm{I}_{p,q}$ \\
\midrule
\multirow{2}{*}{$\mathfrak{c}_n$} & $\Sp(2n,\R)$ & $\sigma(X)=\bar{X}$ & $\theta(X)=\mathrm{J}_n X\mathrm{J}_n^{-1}$ \\
 & $\Sp(2p,2q)$ & $\sigma(X)=-\mathrm{K}_{p,q}X^*\mathrm{K}_{p,q}.$ & $\theta(X)=\mathrm{K}_{p,q}X\mathrm{K}_{p,q}$ \\
\midrule
\multirow{2}{*}{$\mathfrak{d}_n$} & $\SO(p,q)$ & $\sigma(X)=\mathrm{I}_{p,q}\bar{X}\mathrm{I}_{p,q}.$ & $\theta(X)=\mathrm{I}_{p,q}X\mathrm{I}_{p,q}$ \\
 & $\SO^*(2m)$ & $\sigma(X)=\mathrm{J}_m\bar{X}\mathrm{J}_m^{-1}.$ & $\theta(X)=\mathrm{J}_m X\mathrm{J}_m^{-1}$ \\
\bottomrule
\end{tabular}
\tabtitle{4.8}{经典李代数 $G_\C$ 的非紧形式 $G$}
\end{table}

% [Tutor Note: p.31 表 4.8 可用 $\theta=\rho\sigma$ 验算：以 $\mathfrak{a}_n$ 首行为例，由表 4.7 有 $\rho(X)=-\bar{X}^t$，故 $\theta(X)=\rho(\sigma(X))=-\overline{\sigma(X)}^t=-\overline{\bar{X}}^t=-X^t$，与表中一致；$\mathfrak{b}_n$ 与 $\mathfrak{d}_n$ 的 $\SO(p,q)$ 两行公式完全相同，是因为它们分别是 $\mathfrak{so}(2n+1)$ 与 $\mathfrak{so}(2n)$ 的不同实现。]

\begin{exmp}{4.9}
可以把 $\SO(p,p+q)$-Higgs 丛看成 $\SO(2p+q,\C)$ 的模空间上对合
\[
\Theta_{\SO(p,p+q)}\colon (E,\Phi)\mapsto(E,-\Phi)
\]
的不动点；这些不动点对应这样的向量丛 $E$：$E$ 有一个共轭于 $\mathrm{I}_{p,p+q}$ 的自同构 $f$，它把 $\Phi$ 送到 $-\Phi$，并且 $E$ 的 $\pm1$ 特征空间的维数分别为 $p$ 与 $p+q$。
\end{exmp}

% [Tutor Note: p.31 例 4.9 的维数核对：$f$ 的 $+1$ 与 $-1$ 特征空间维数之和 $p+(p+q)=2p+q$，恰为 $\SO(2p+q,\C)$ 标准表示的维数，与 "$f$ 共轭于 $\mathrm{I}_{p,p+q}$" 一致。]

下面我们将概述 G-Higgs 丛的 $(\mathrm{B},\mathrm{A},\mathrm{A})$-膜作为式 (4.5) 中 $i_1$ 在 Hitchin 纤维化上诱导作用的不动点集的各种不同表现。特别地，这些膜与 Hitchin 纤维化的正则纤维可以以非常不同的方式相交：

\begin{itemize}[leftmargin=1.4em,itemsep=1pt,topsep=1pt]
\item \textbf{与正则纤维的有限交。}如前所述，由 [89, Theorem 4.12]，当 G 是 $G_\C$ 的分裂实形式时，G-Higgs 丛的 $(\mathrm{B},\mathrm{A},\mathrm{A})$-膜与 Hitchin 纤维化的纤维相交于二阶点。由于我们在前面几章已经概述过这一情形，这里不再进一步讨论（关于分裂实形式的谱数据，见 [11, 29, 88, 89, 92, 97] 等文献）。
\item \textbf{与正则纤维的正维交。}$\U(p,p)$-Higgs 丛在 $\cM_{\Gl(2p,\C)}$ 中的 $(\mathrm{B},\mathrm{A},\mathrm{A})$-膜（如习题 \textbf{II.(4)} 所示），以及 $\SU(p,p)$-Higgs 丛在 $\cM_{\Sl(2p,\C)}$ 中的 $(\mathrm{B},\mathrm{A},\mathrm{A})$-膜，与正则纤维相交于由雅可比簇或商谱曲线的 Prym 簇给出的有限维阿贝尔簇，如 [89, 91] 所示。
\end{itemize}
